\section{Are DLP fragments still worthwhile alongside OWL 2 RL?}
\label{sec:profiles}
Our answer separates theoretical and computational merit from practical
adoption. A fragment can be worthwhile even when its prevalence in deployed
ontologies is unknown. DLP remains useful as a foundation and as a decomposition
of evaluator capabilities. OWL 2 RL explicitly credits DLP among its
influences~\cite{owlprofiles2012}; standardization incorporated the central
idea rather than making it irrelevant.

In particular, DLP L3's existential consequents can express useful statements
beyond RL's syntax. If specialized DLP L3 reasoning is faster than a more expressive
OWL reasoner on the same DLP L3 inputs and reasoning tasks, while preserving the
same correctness and completeness obligations, that is evidence for a useful
computational fragment. No estimate of how often users need only DLP L3 is required
for this conclusion. The finite existential-bridge study in
Section~\ref{sec:hermit} now supplies such scoped computational evidence:
the preserved baseline and candidate both answer the same consistency and
named-instance task faster than the pinned broader-language HermiT
distribution~\cite{glimm2014hermit}. All 27 measured workers match independently
justified answers. The candidate/HermiT paired median task ratios are
0.046429, 0.262592, and 0.378055 for 100, 1,000, and 10,000 named A instances;
the corresponding full-process ratios are 0.379001, 0.475576, and 0.415056.
These are results for one controlled family and three process blocks, not a
new complexity bound or a comparison with every broader-language OWL implementation.
Most of this advantage is already present in the fixed baseline.

Measured advantages and mathematical guarantees provide different evidence.
One can prove an expressiveness--complexity tradeoff without an adoption survey;
one can also demonstrate an implementation advantage without proving a new
complexity separation. Hustadt, Motik, and Sattler~\cite{hustadt2005} establish
P-complete data complexity for basic Horn-SHIQ reasoning, while general SHIQ
instance checking is coNP-complete under their assumptions. This is an example
of prior theoretical justification for expressive Horn fragments. It does not
identify our DLP L3 with Horn-SHIQ or give our unrestricted witness chase that bound.
Termination, combined complexity, data complexity, and query complexity require
their own precise language and task definitions.
In particular, nontermination of this Skolem chase is a limitation of the
chosen procedure, not by itself a proof that reasoning in the original
description-logic fragment is undecidable.

Three classifications must remain distinct: the thesis's constructor-position
fragments, the actual RDF encodings accepted by this compiler, and the
conditions under which an optimization is valid. An L0 ontology can contain
expensive transitive binary joins. An ontology accepted with the L2 ceiling may
contain only unary implications. Choosing a higher ceiling does not add absent
equality or constraints, and a label alone cannot predict cost.

\subsection{Overlap rather than automatic profile inclusion}
Concrete syntax witnesses rule out identifying these labels with RL subsets.
The DLP L3 inclusion $A\sqsubseteq\exists R.B$ has a general existential
consequent absent from RL's superclass grammar. Even L1's $\top\sqsubseteq A$
is outside its subclass grammar. Conversely, RL supports restricted property
chains and keys, which this compiler rejects~\cite{owlprofiles2012}. These
are statements about syntactic membership, not impossibility of equivalent
rewriting or a claim that DLP L3 contains RL.

Global restrictions matter independently. The compiler accepts simultaneous
transitivity and functionality of the same property. That rule set has a
first-order interpretation, but the property fails OWL 2 DL's simplicity
condition for functionality~\cite{owlsyntax2012}. Parser acceptance therefore
does not certify a standard OWL 2 RL ontology. Datatype completeness is another
separate gap in this implementation.

OWL 2 RL/RDF rules can also be applied beyond syntactically conforming RL.
Their Direct-Semantics correspondence has explicit conditions, including
restrictions on entity roles and on the conclusions covered by the theorem
\cite{owlprofiles2012,owlconformance2012}. A comparison must choose its
semantic contract. Our equality acts on individual arguments, not predicate
identifiers. Requiring identity with every triple produced by an RDF-Based
reasoner can therefore demand consequences outside the intended contract.

\subsection{What restrictions can still buy}
For this implementation, function freedom gives finite active-domain
materialization. Eligible states without explicit or datatype-induced equality
avoid representative splitting during
deletion. Same-variable unary bodies permit set-intersection evaluation.
These properties are useful internal dispatch conditions, including within a
broader supported ontology. They do not require presenting another public
ontology language as a replacement for RL.

A small compiler can also be easier to audit and reproduce. Its value depends
on real schemas fitting the contract and on explicit rejection of unsupported
axioms. Conversion costs, excluded standard features, and additional
conformance work are part of the cost of the restriction. The earlier
certificate and unary experiments demonstrate workload-specific benefits and
failures, not an advantage of requiring a DLP frontend over accepting the same
axioms through a conforming RL frontend.

LUBM illustrates the distinction between ontology coverage and query
sufficiency. Its full ontology needs this implementation's DLP L3 mode, and all
fourteen named answer sets match. That does not prove those questions require
generated witnesses. Nor would matching them after a projection prove ontology
equivalence. Existential consequences can affect named answers: from $A(a)$,
$A\sqsubseteq\exists R.B$, and $\exists R.B\sqsubseteq C$, we obtain $C(a)$
\cite{owldirect2012}. Dropping the witness-producing inclusion without a
justified rewriting loses this consequence.
The HermiT study evaluates precisely this bridge together with a B-only
negative individual and an unrelated empty-answer class. Both GCIs are
essential to the answer, but explicit witness construction is not: the TBox
entails $A\sqsubseteq C$, and HermiT retains its normal optimizations. The
family also belongs to OWL 2 EL, so the result does not establish superiority
over specialized EL reasoners. It compares the specified entailment task,
not entire materializations or generated witness identities.

\subsection{Practical adoption is a separate question}
OWL 2 EL and QL supply other established tradeoffs, respectively oriented toward
large terminologies with existential descriptions and query rewriting over
large instance data~\cite{owlprimer2012}. They are alternatives to evaluate,
not automatic solutions for every DLP L3 query. In particular, avoiding explicit
infinite witness expansion for classification does not establish completeness
for arbitrary conjunctive queries or cardinality constraints.

The frequency with which real applications need only DLP L3 is a difficult empirical
question. Published ontology collections can reveal supported constructs in
those collections, but may be unrepresentative; they also do not fully reveal
application queries, update needs, or equivalent rewritings. We do not require
an answer to that question before recognizing a fragment's theoretical value.

For practical interoperability claims, checked standard profiles and explicit
extensions would complement the historical DLP L0--L3 interface. A deployment study
could measure coverage and rejection reasons, full input-to-answer and update
costs, conversion and memory, and ablations with compiled rules held fixed.
Its coverage conclusions would remain specific to the sampled population.
Current evidence supports the fragments as foundations and research targets
and demonstrates a computational advantage on the stated existential task.
That scoped result stands independently of prevalence; it does not establish
the historical interface as a better interoperability standard or settle
performance on other ontology and query families.
